\chapter{Raisonner, définir et démontrer}
\section*{Exercice 1.1 -- Nier sans ambiguïté}
Nier : « tous les termes de la suite sont positifs » ; « il existe un entier $n$ tel que $u_n>3$ ». Préciser le sens choisi pour « positif ».

\section*{Exercice 1.2 -- Implication et réciproque}
Démontrer : si $n$ est un entier impair, $n^2$ est impair. La réciproque est-elle vraie ?

\section*{Exercice 1.3 -- Images réciproques}
Pour une fonction $f:E\to F$ et $A,B\subset F$, prouver $f^{-1}(A\cap B)=f^{-1}(A)\cap f^{-1}(B)$.

\section*{Exercice 1.4 -- Compter une réunion}
Dans une classe de $32$ élèves, $18$ pratiquent A, $15$ pratiquent B et $7$ les deux. Combien pratiquent au moins une activité ? Aucune ?

\section*{Exercice 1.5 -- Une récurrence facultative}
Montrer par récurrence que $2^n\ge n+1$ pour tout $n\in\N$.

\section*{Exercice 1.6 -- Ordre des quantificateurs}
Comparer : « pour tout réel $x$, il existe un réel $y$ tel que $y>x$ » et « il existe un réel $y$ tel que pour tout réel $x$, $y>x$ ».
